\section*{Resumen}

La caracterización del terreno sustenta decisiones estructurales y de control de vital importancia durante el proyecto, la ejecución y el seguimiento de una obra civil. Este trabajo desarrolla una plataforma electrónica orientada al análisis multicanal de ondas superficiales (MASW), compuesta por nodos de captación con geófonos SM-24, acondicionamiento y adquisición en PSoC~5LP, coordinación inalámbrica mediante ESP32 y un nodo instrumentado para registrar la excitación. Las evidencias obtenidas en campo permitieron construir un registro espacial con un receptor desplazado y reconocer, mediante análisis espectral, una relación señal--ruido insuficiente en la banda baja de interés. Estos resultados determinan los siguientes hitos: aumentar la energía y repetibilidad de la excitación, completar la validación temporal entre nodos y extender el análisis a la formación de imágenes de dispersión, selección modal e inversión.

\section{Introducción}
\label{sec:introduccion}

En una obra civil, conocer el terreno condiciona decisiones estructurales y de control que son de vital importancia para su seguridad y desempeño. Antes de construir se debe reconocer la sucesión de materiales y estimar parámetros con los cuales dimensionar la fundación. Durante la ejecución se controla la uniformidad de rellenos y capas compactadas; posteriormente, nuevas mediciones permiten investigar cambios, daños o sectores que requieren mantenimiento. El reconocimiento inicial constituye un requisito de diseño, mientras que el seguimiento posterior exige observar el terreno sin alterar la estructura existente \citep{FHWA2006,Park2018Roads}.

Los ensayos geotécnicos tradicionales aportan información directa y localizada. El ensayo de penetración estándar (SPT) combina resistencia a la penetración con recuperación de muestras en puntos discretos, mientras que CPT y CPTu generan perfiles continuos durante la penetración. Los ensayos \emph{downhole} y \emph{crosshole} permiten estimar propiedades dinámicas entre puntos conocidos, pero requieren uno o varios pozos. Por su naturaleza, estas técnicas resultan indispensables cuando se necesitan muestras o verificaciones puntuales, aunque su intervención física y cobertura localizada dificultan extenderlas a toda la obra o repetirlas después de construir \citep{FHWA2006,Foti2014}.

Como complemento, los procedimientos de superficie aplican una perturbación y observan variables medibles desde la misma superficie. Las características temporales, espaciales y espectrales de la respuesta contienen información sobre las propiedades mecánicas del medio. De este modo puede realizarse un reconocimiento preliminar antes de perforar, localizar con mejor criterio los ensayos invasivos, completar la información entre sondeos y repetir la observación durante la vida de la obra. Ambos enfoques son, por tanto, complementarios: el reconocimiento superficial orienta dónde investigar con mayor detalle y los ensayos directos aportan la comprobación localizada \citep{Foti2014,Foti2018}.

Los métodos de superficie se han popularizado por estas ventajas, pero su acceso sigue siendo difícil para muchos interesados. Un arreglo requiere sensores de costo elevado, instrumentación analógica delicada, conversión, almacenamiento, coordinación temporal y procesamiento de las señales. El propósito de este trabajo es afrontar en conjunto esos problemas y desarrollar una plataforma cómoda de operar, abierta a ampliaciones y de menor costo relativo, sin perder las magnitudes necesarias para interpretar los registros.

Se propone utilizar el PSoC~5LP por su versatilidad y por la combinación poco habitual, entre microcontroladores comerciales, de bloques analógicos configurables, conversión y lógica digital en un mismo dispositivo. En conjunto con un ESP32, cada nodo puede acondicionar, adquirir, almacenar, coordinarse inalámbricamente y presentar su estado al operador. También se instrumenta el elemento de excitación y se lo incorpora como un nodo más de la red, de manera que cada respuesta quede asociada con el impacto que la produjo.

Si bien la disponibilidad inicial de un solo receptor obliga a repetir la adquisición en posiciones sucesivas, el diseño permanece abierto y fue concebido para incorporar varios nodos receptores. Para establecer qué deben medir y preservar esos nodos, primero es necesario comprender la respuesta mecánica del terreno, las formas de propagación que admite y el tipo de movimiento que puede observarse desde la superficie.
